\documentclass[12pt, french, latin.medieval, a4paper]{report}
% package de langues 
\usepackage[french]{babel}
\usepackage[utf8]{inputenc}
\usepackage[T1]{fontenc}

% package pour référence de liens sur pdf 
\usepackage[pdftex]{hyperref} % liens
\usepackage{lastpage} % numéro de page
\usepackage{caption}

% package de formes de document 
\usepackage[left=2.5cm,
              right=2.5cm,
              top=2.5cm,
              bottom=2.5cm
             ]{geometry}
% pour les tableaux 
\usepackage{array,multirow,makecell} 
\newcolumntype{R}[1]{>{\raggedleft\arraybackslash }b{#1}}
\newcolumntype{L}[1]{>{\raggedright\arraybackslash }b{#1}}
\newcolumntype{C}[1]{>{\centering\arraybackslash }b{#1}}

% package pour la page de couverture 
\usepackage{graphicx}
\usepackage[cyr]{aeguill}
\usepackage{fancyheadings}
\usepackage{amsmath}
\usepackage{amsfonts}
\usepackage{amssymb}
\usepackage{geometry}
\usepackage{eso-pic}
\usepackage{tikz}
\usetikzlibrary{calc}


% nouvelle commande pour la forme 
%\fancyhf[OFC]{\thepage / \pageref{LastPage}}
\rfoot{\thepage / \pageref{LastPage}}
%\setcounter{tocdepth}{3}

% pakages de couleurs du document
\usepackage{color}
\usepackage{xcolor}


% package pour la chimie
\usepackage[version=4]{mhchem}


% package pour la visualisation du code
\usepackage{listings}
\definecolor{darkWhite}{rgb}{0.94,0.94,0.94}
 
\lstset{
  aboveskip=3mm,
  belowskip=-2mm,
  backgroundcolor=\color{darkWhite},
  basicstyle=\footnotesize,
  breakatwhitespace=false,
  breaklines=true,
  captionpos=b,
  commentstyle=\color{red},
  deletekeywords={...},
  escapeinside={\%*}{*)},
  extendedchars=true,
  framexleftmargin=16pt,
  framextopmargin=3pt,
  framexbottommargin=6pt,
  frame=tb,
  keepspaces=true,
  keywordstyle=\color{blue},
  language=python,
  literate=
  {²}{{\textsuperscript{2}}}1
  {⁴}{{\textsuperscript{4}}}1
  {⁶}{{\textsuperscript{6}}}1
  {⁸}{{\textsuperscript{8}}}1
  {€}{{\euro{}}}1
  {é}{{\'e}}1
  {è}{{\`{e}}}1
  {ê}{{\^{e}}}1
  {ë}{{\¨{e}}}1
  {É}{{\'{E}}}1
  {Ê}{{\^{E}}}1
  {û}{{\^{u}}}1
  {ù}{{\`{u}}}1
  {â}{{\^{a}}}1
  {à}{{\`{a}}}1
  {á}{{\'{a}}}1
  {ã}{{\~{a}}}1
  {Á}{{\'{A}}}1
  {Â}{{\^{A}}}1
  {Ã}{{\~{A}}}1
  {ç}{{\c{c}}}1
  {Ç}{{\c{C}}}1
  {õ}{{\~{o}}}1
  {ó}{{\'{o}}}1
  {ô}{{\^{o}}}1
  {Õ}{{\~{O}}}1
  {Ó}{{\'{O}}}1
  {Ô}{{\^{O}}}1
  {î}{{\^{i}}}1
  {Î}{{\^{I}}}1
  {í}{{\'{i}}}1
  {Í}{{\~{Í}}}1,
  morekeywords={*,...},
  numbers=left,
  numbersep=10pt,
  numberstyle=\tiny\color{black},
  rulecolor=\color{black},
  showspaces=false,
  showstringspaces=false,
  showtabs=false,
  stepnumber=1,
  stringstyle=\color{gray},
  tabsize=4,
  title=\lstname,
}

%package pour les algorithmes
\usepackage{algorithm}
\usepackage{algorithmic}
\usepackage{amsmath}

% package pour les tableaux
\usepackage{longtable}
\usepackage{booktabs}

% package pour les figures 
\usepackage{subcaption}

% package de reférencement
% \usepackage[backend=biber,safeinputenc, sorting=none]{biblatex}

% mettre la bibliographie dans la table de matières
%\usepackage[nottoc, notlof, notlot]{tocbibind}
\include{page_couverture.tex}


\author{MBE MBE MINDJANA LOIC HENRI}
\title{GNPS et MDByTE}

\begin{document}
    
    \pagecouverture
    {Multi-view Clustering : A Survey}
    {Yan Yang et Hao Wang}
    \tableofcontents
    \vspace{20cm}
    
    \section{Introduction}
    De plus en plus de nos jours les données sont obtenues de diverses sources qui sont généralement liées à des taches différentes sur des objets communs. Par exemple les données sur les voitures qu'on retrouvent généralemet sur des sites sont leurs images et descriptions textuelles qui constituent des représentations différentes (\textbf{des vues}) de ces derniéres. Afin de mieux explorer ce type de données, on utilise des algorithmes de régroupement spécifiques appellés \textbf{Algorithmes de Clustering Multi-vues} qui permèttent de tenir compte de la complémentarité et du consensus (accord) entre les différentes perspectives des données communément appélées \textbf{vues}. une enquète a été ménée par les chercheurs \textbf{Yan Yang} et \textbf{Hao Wang} pour recenser l'ensemble de ces algorithmes. Ce recensement a conduit à l'article \textbf{Multi-view Clustering: A Survey} de référence \cite{yang2018multi} qui sera le sujet de ce document. Dans les lignes qui suivent nous presenterons d'abord l'article ensuite nous parlerons des algorithmes de clustering Multi-vues telque décrit dans l'article et enfin nous listerons des datasets y présentées.
    
    \section{Description de l'article}
    L'article \textbf{Multi-view Clustering: A Survey} \cite{yang2018multi} est le résultat d'une enquète ménée sur les algorithmes de clustering multi-vues fait par les chercheurs Yan Yang et Hao Wang. Dans cet article, les auteurs décrivent les principes globaux des algorithmes de clustering et donne une topologie de ces derniers. Les informations liées à l'article sont les suivantes:
    \begin{itemize}
        \item \textbf{Titre}: Multi-view Clustering: A Survey
        \item \textbf{Auteurs}: Yan Yang et Hao Wang
        \item \textbf{Année de publication}: 2018
        \item \textbf{Journal}: Big Data Mining and Analytics
    \end{itemize}
    
    \section{Clustering multi-vues}
    Le clustering multi-vues est une technique d'apprentissage non supervisée qui consiste à regrouper des données(observations) similaires entre eux et d'en éloigner ceux qui ne le sont pas en tenant compte qu'elles sont constituées de différentes vues. Les vues sont perçues comme les différentes représentations des observations utilisées. 
    
    Le clustering multi-vues se différentie du "clustering" principalement par le fait qu'on cherche à exploiter la complémentarité et l'accord (consensus) entre des vues de données dans le but de mieux les regrouper et qu'on a qu'une seule vue en "clustering". La complémentarité et le consensus entre les vues sont donc des notions importantes à définir:
    \begin{itemize}
        \item \textbf{Complémentarité} : la complémentarité entre deux vues V1 et V2 est décrit par le fait qu'il existe sur des observations, des informations qui ne peuvent être représentées que par V1 et d'autres que par V2
        \item \textbf{Consensus}: le consensus entre deux vues V1 et V2 est décrit par le fait qu'il existe sur des observations, des informations qui peuvent être décrite par V1 et par V2
    \end{itemize}
    
    De là on a deux principes qui en ressorte pour mieux faire de l'apprentissage de façon globale:
    \begin{itemize}
        \item \textbf{Principe associé à la complémentarité}: les vues doivent être utilisées afin de décrire les observations de manière plus complète et plus précise, vu que chaque vue ressort des tâches spécifiques qui apportent des informations complémentaires
        \item \textbf{Principe associé au consensus}: il faut maximiser le consensus ou encore la cohérence entre les vues des observations afin de mieux cerner les informations complémentaires entre les vues. La raison de la maximisation vient de l'étude des chercheurs Daguplaet al. qui prouvent dans l'article de référence \cite{dasgupta2001pac} que la probabilité d'erreur associée au désaccord entre les vues est supérieur à la somme de ceux des vues séparément.
    \end{itemize}

    Ces deux principes constituent la base sur laquelle les intuitions des algorithmes de clustering multi-vues s'appuient. Des manières de les exploiter, on conduit à l'obtention des types d'algorithmes qui seront présentés dans les sous sections suivantes.
    
        \subsection{Algorithmes de co-formation}
         Dans cette catégorie d'algorithme, on cherche à maximiser l'accord entre les vues tout en améliorant les clusters itérativement de vues en vues. On construit initialement les clusters pour chaque vue, et les clusters d'une vue sont améliorés en utilisant ceux d'autres vues dans le but de maximiser le consensus entre eux. On procède ainsi itérativement tour à tour pour chaque vue jusqu'à obtenir une convergence du consensus entre les vues ou après un certain nombre d'itérations.
        
        La figure \ref{fig:co-train1} l'illustre assez bien. On construit des clusters pour chacune des vues 1 et 2 après entraînement à la suite on améliore les clusters d'une vue, en utilisant des connaissances de l'autre dans le but maximiser la cohésion entre les vues.
        
        \begin{figure}
            \centering
            \includegraphics[width=\linewidth]{Images/co-training.png}
            \caption{Processus général des algorithmes de co-formation \cite{yang2018multi}}
            \label{fig:co-train1}
        \end{figure}
        On peut dès lors comprendre que ces algorithmes possèdent un défaut majeur qui est le temps d'exécution, car on reconstruit les clusters en un certain nombre d'itérations.

        \subsection{Algorithmes multi-noyaux}
        S'il est difficile de construire des clusters dans un espace d'une dimension donnée, on pourrait plus aisément le faire en projetant les observations dans un espace correspondant de plus grande dimension. C'est sur ce principe que les algorithmes de clustering par noyaux sont construits pour les données représentées par des caractéristiques d'une même représentation. On recherche donc une fonction noyau qui permettra d'obtenir de meilleurs clusters suivant un critère d'évaluation comme l'inertie intra-classe.
        
        Pour cette catégorie de clustering multi-vues, on cherche également une fonction noyau qui permettra d'obtenir des clusters dans un espace de plus grande dimension qui agglomère les espaces des différentes vues. Ceci est fait en combinant linéairement ou non linéairement des fonctions noyaux obtenues, sur chacunes des vues dans le but d'obtenir des clusters en tenant compte le plus possible des principes de complémentarité et de consensus donnés plus haut. La figure \ref{fig:multi-kernel1} décrit assez bien le processus global des algorithmes de ce type, on a des vues $View_1,View_2,...,View_m$, dont on obtiendra des fonctions noyaux correspondantes qui seront combinées linéairement ou non linéairement par apprentissage pour obtenir un noyau unifié.

        Généralement ces algorithmes s'inspirent d'autres algorithmes de clustering à une vue tel que KMEANS, et les modifient pour intégrer les fonctions noyaux dans le contexte des données multivues. Dans le cas de Kmeans, on définit la fonction noyau unifiée $K$ comme dans l'équation \ref{equa:kernel1}. vu qu'on souhaite minimiser l'inertie intraclasse du clustering $I$, on aura le problème d'optimisation \ref{equa:kernel2} où $W$ est la matrice de poids permettant d'exploiter les deux principes de consensus et de complémentarité. Il faut noter que $\gamma$ est la matrice indicateur de clusters c'est à dire que si l'observation $i$ est dans le cluster $k$ alors $\gamma_{ik}=1$ sinon $\gamma_{ik}=0$.
        \begin{equation}
            K(x,y) = {[\exp(-\frac{{||x-y||}^2}{2 \sigma^2}) + R]}^P
            \label{equa:kernel1}
        \end{equation}

        \begin{equation}
            \begin{aligned}
                \min_{W,R,P} I = \sum_{v=1}^{V} \sum_{i=1}^{n} w_{i,v} \sum_{k=1}^{K} \gamma_{ik} {|| \phi(x_{i}^{v}) - \phi(c_{j}^{v}) ||}^2 
                \\
                = \sum_{v=1}^{V} \sum_{i=1}^{n} w_{i,v} \sum_{k=1}^{K} 2 \gamma_{ik} [{(1 + R)}^P + K(x_{i}^{v}, c_{j}^{v})]
                \\ 
                \text{s.t, } w_{i,v} \geq 0\text{, }\prod_{v} w_{i,v}=1
            \end{aligned}
            \label{equa:kernel2}
        \end{equation}
        
        \begin{figure}
            \centering
            \includegraphics[width = \linewidth]{Images/multi-kernel.png}
            \caption{Processus général des algorithmes multi-noyaux \cite{yang2018multi}}
            \label{fig:multi-kernel1}
        \end{figure}

        \subsection{Algorithmes basés sur les graphes}
        Pour un ensemble d'observations n'ayant qu'une vue, on peut construire un graphe de similarité(très proche d'un réseau moléculaire) qui permettrait après découpe (clustering spectral) d'obtenir des clusters.
        
        L'idée \ref{fig:graph1} pour pouvoir tenir compte des données à plusieurs vues, est de construire pour chaque vue des graphes de similarités. Puis de définir une politique de fusion de graphe qui s'appuie généralement sur l'attribution des poids à des observations, par rapport aux groupes formés dans le but de maximiser une fonction objective. De là on obtient un graphe unifié de similarité qui pourra être sujet d'un clustering spectral pour pouvoir construire des clusters.

        \begin{figure}
            \centering
            \includegraphics[width=\linewidth]{Images/graph.png}
            \caption{Processus général des algorithmes basés sur les graphes \cite{yang2018multi}}
            \label{fig:graph1}
        \end{figure}
        
        \subsection{Algorithmes basés sur les réseaux}
        
        Les algorithmes basés sur les graphes ne permettent pas de tenir compte des relations entre les observations des différentes vues. Dans certains cas, une observation dans une vue peut avoir plusieurs équivalents dans d'autres vues. Pour tenir compte de ces relations, on construit des graphes pour chaque vue qui seront désormais considérées comme des réseaux, puis on construira un réseau global de ces réseaux \ref{fig:network1}. Ainsi par optimisation d'une fonction incluant les principes de consensus et de complémentarité, les indicateurs de clusters de chacun des vues seront obtenus.
    
        \begin{figure}
            \centering
            \includegraphics[width =\linewidth]{Images/netwrk.png}
            \caption{Réseaux de réseaux \cite{yang2018multi}}
            \label{fig:network1}
        \end{figure}

        \subsection{Algorithmes spectraux}
        Effectuer du clustering revient à trouver k groupes d'observations similaires entre eux. Cela se fait par une fonction $f$ qui permet d'attribuer à chaque observation un indicateur propre à un des $K$ clusters. On peut donc voir que $f$ est lui même un indicateur pour attribuer une observation à un cluster donné et donc peut être vu comme une matrice $(n \times k)$ où $f_{i,k}$ est la probabilité que l'observation $i$ soit dans la classe $k$.
        
        L'objectif des algorithmes de clustering de ce type est de calculer cette matrice tout en tenant compte des similarités entre les observations(par ricochet d'un graphe de similarité), en cherchant à résoudre le problème d'optimisation \ref{equa:spectral1} où $L$ est le laplacien du graphe de similarité des observations.
        
        \begin{equation}
            \min_{F^{T}F = I} tr(F^{T}LF)
            \label{equa:spectral1}
        \end{equation}

        Pour pouvoir prendre en compte plusieurs vues, le problème d'optimisation \ref{equa:spectral1} change pour le problème \ref{equa:spectral2} qui est équivalent au problème \ref{equa:spectral3} où $L = \sum_{v=1}^{V} \alpha_{v}^{r} \times L_{v}$, avec $r$ étant un coefficient pour contrôler les distributions des vues. Nous remarquons donc que ce problème est à peu près identique au problème de clustering spectral à une vue.

        \begin{equation}
            \begin{aligned}
                \min_{F^{T}F = I, \alpha} \sum_{v=1}^{V} \alpha_{v}^{r} tr(F^{T}L_{v}F) \\
                \text{s.t, } \sum_{v} \alpha_{v} = 1 \text{ , } \alpha_{v} \geq 0
            \end{aligned}
            \label{equa:spectral2}
        \end{equation}

        \begin{equation}
            \begin{aligned}
                \min_{F^{T}F = I, \alpha} tr(F^{T}LF) \\
                \text{s.t, } \sum_{v} \alpha_{v} = 1 \text{ , } \alpha_{v} \geq 0
            \end{aligned}
            \label{equa:spectral3}
        \end{equation}

        Les vecteurs propres dans l'ordre croissant de leurs valeurs propres du laplacien conjugué $L$ seront calculés, puis soumis à un algorithme de clustering classique pour pouvoir trouver $F$. Il faut noter qu'avec cette méthode de clustering, il est possible de détecter le nombre de clusters approprié (qui est le nombre de vecteurs propres), en s'appuyant sur les écarts entre les valeurs propres consécutives. Le nombre $K$ de clusters peut être choisi ou obtenu en détectant une grande variation entre deux valeurs propres consécutives selon l'ordre croissant de ces valeurs.

        \subsection{Algorithmes de sous-espaces}
        
        L'idée ici est de rechercher une représentation unifiée des vues pour ensuite y effectuer un algorithme de clustering à une vue. De ce fait on procède généralement de deux manières tel que sur la figure \ref{fig:subspace}:
        \begin{enumerate}
            \item soit on apprend directement la représentaion depuis les vues des données

            On cherche à trouver généralement un ensemble de matrices $W_{v}(d_{v}, d')$ de poids positives attachées aux différentes vues et une matrice $H(n \times d')$ qui sera la représentation unifiée des données par rapport aux vues. Ici $d_{v}$ est la dimension des données suivant la vue $v$, $d'$ la dimension des données de la représentation unifiée et $n$ est le nombre d'observations. Le problème d'optimisation générale utilisé par les algorithmes de ce type est celui de l'équation \ref{equa:subspace1}, ce problème sera modifié en fonction des spécificités. Il faut noter que la norme $||A||_{F}^{2} = tr(A.A_{T}) = \sum_{i,j} a_{i,j}$ est appelée la norme de Frobenius.
            
            \item soit on apprend d'abord un espace latent avant de l'utiliser pour apprendre la représentation unifiée: généralement on effectue une projection des données des vues dans l'espace latent en question           
        \end{enumerate}

        \begin{figure}
            \centering
            \includegraphics[width = \linewidth]{Images/subspace.png}
            \caption{Processus général des algorithmes de sous-espaces \cite{yang2018multi}}
            \label{fig:subspace}
        \end{figure}

        \begin{equation}
            \begin{aligned}
                \min_{W,H}  {|| X - WH ||}_{F}^{2} \\
                \text{s.t, } W \geq 0 \text{, }H \geq 0
            \end{aligned}
            \label{equa:subspace1}
        \end{equation}

    \section{Datasets recensés}
    Les auteurs au cours de leur enquête ont recensé un ensemble de bases de données d'entraînement multi-vues, qui nous ont interpellés. Elles sont les suivantes:
    \begin{itemize}
        \item \href{http://mlg.ucd.ie/datasets/bbc.html}{3Sources Datasets}: est un dépôt qui contient des corpus de textes multivues d'articles de 3 journaux, des datasets multivues de Twitter et bien d'autres
        \item \href{ttp://www.cs.cmu.edu/afs/cs.cmu.edu/project/theo-20/www/data/}{WebKb Datasets} : ensemble de datasets contenant des données de pages web issues de différentes université
        \item \href{http://lig-membres.imag.fr/grimal/data.html}{Newsgroup Datasets}: est un dépôt contenant des datasets multivues tels que "Cora Dataset", Reuters Multilingual Dataset.

        \item \href{http://archive.ics.uci.edu/ml/datasets/Multiple+Features}{Handwritten Digit Dataset} cette dataset contient comme vues les différents types de caractéristiques qu'on peut extraire des images des chiffres

        \item \href{https://archive.ics.uci.edu/ml/datasets/One-hundred+plant+species+leaves+data+set}{100leaves Dataset} Ici on a des données multivues sur des feuilles de plantes.

        \item \href{https://archive.ics.uci.edu/ml/datasets/YouTube+Multiview+Video+Games+Dataset}{YouTube Video Dataset}: cette dataset contient différents types de caractéristiques de vidéos youtube
    \end{itemize}

    \section{Difficultés}
    Lors de la lecture de l'article, la principale difficulté a été le fait que l'article contenait beaucoup de nouvelles notions qui fallait élucider afin de comprendre les types d'algorithmes. Actuellement nous ne comprenons pas encore les algorithmes de clustering multi-vues multi-tâches, en raison de leur absence dans le rapport.

    \section{Conclusion}
    En conclusion nous avons parlé de la topologie des algorithmes de clustering multi-vues donnée par les auteurs et avons listé des datasets qu'ils ont recensé et qui nous ont particulièrement interpellé. Les différents types d'algorithmes suivent des intuitions de base différentes, d'autres sont plus utiles pour un certain type de données et d'autres non. Au vue cette topologie, les algorithmes de clustering multi-vues de sous espaces, spectraux et basés sur les graphes ont attirés notre attention, du fait de leurs principes de base qui s'attèlent mieux aux données MS et RMN, et notre logique. Nous comptons dans les jours à venir à plus étudier ces types d'algorithmes respectivement.
   
\bibliographystyle{plain}
    \bibliography{reference_GNPS_MADByTE}
\end{document}